\pagebreak
\section{Метод решения}
Корпус --- 1\,000\,000 документов из ЛР1 (34\,818 статей Wikivoyage, 620\,164 статьи Wikipedia, 48\,768 вопросов Travel Stack Exchange с ответами, 296\,250 постов Reddit; тексты \texttt{lab1/data\_1m/text}, 5,6~ГиБ; заголовки и ссылки --- в \texttt{lab1/data\_1m/raw/<источник>/meta.jsonl}).

Программы:
\begin{itemize}
    \item \texttt{scripts/doclist.py} (Python) --- список документов \texttt{data/docs.tsv}: путь к тексту, заголовок, URL через табуляцию. Номер документа (docID) --- номер строки с нуля: Wikivoyage --- 0--34\,817, Wikipedia --- 34\,818--654\,981, Travel SE --- 654\,982--703\,749, Reddit --- 703\,750--999\,999. Источники скрипт находит сам по каталогам корпуса.
    \item \texttt{build/indexer} (C++) --- индексатор: \texttt{indexer data/index.bin < data/docs.tsv}, печатает статистику и время.
    \item \texttt{build/lookup} (C++) --- чтение индекса: \texttt{lookup index.bin} выводит для каждого слова со стандартного ввода число документов (df) и их номера; \texttt{lookup index.bin -d} по номерам документов выводит заголовок, URL и число токенов. Утилита проверяет, что индекс читается, и нужна для автотестов; булев поиск --- следующая работа.
\end{itemize}
Структуры данных индекса (динамический массив, хеш-таблица, сортировка) написаны самостоятельно, без STL, на функциях библиотеки C. STL (\texttt{std::string}) есть только в токенизаторе из ЛР3, что разрешено правилами курса.

\subsection{Термы и внутреннее представление документов}
Текст документа делится на токены токенизатором из ЛР3: токен --- последовательность букв и цифр Юникода, регистр понижен. Терм --- различный токен. Стемминга нет.

Документ после токенизации нигде не хранится целиком: токенизатор передаёт токены (строки UTF-8 в нижнем регистре) функции обратного вызова, и каждый токен сразу добавляется в индекс. От документа остаются только docID в списках документов его термов и запись прямого индекса: заголовок, URL и число токенов (пригодится для ранжирования).

\subsection{Построение индекса}
\begin{enumerate}
    \item Словарь --- хеш-таблица с открытой адресацией и линейным пробированием (хеш FNV-1a, размер --- степень двойки, начальный 4096, удваивается при заполнении на 50\,\%). Для каждого терма хранятся: строка (в общем буфере), хеш, df, номер последнего документа и динамический массив docID (ёмкость удваивается).
    \item Документы обрабатываются по порядку. Для токена ищется терм в таблице (новый добавляется); если номер последнего документа терма не равен текущему, docID дописывается в его список. Так документ попадает в список один раз, а списки сразу упорядочены по возрастанию docID.
    \item После всех документов номера термов сортируются по байтам строки терма (порядок байтов UTF-8 совпадает с порядком кодов Юникода).
    \item Индекс записывается в файл одним последовательным проходом через буфер 64~КиБ.
\end{enumerate}

\subsection{Метод сортировки}
Пары \enquote{терм--документ} не сортируются вообще: документы обрабатываются в порядке docID, поэтому каждый список документов заполняется уже упорядоченным. Сортируется только словарь --- 3,97 млн термов --- восходящей сортировкой слиянием (merge sort) массива номеров термов, сравнение --- \texttt{memcmp} строк.
\begin{itemize}
    \item Достоинства: всегда $O(V \log V)$, без худшего случая быстрой сортировки; слияние читает и пишет массивы последовательно; та же операция слияния --- основа внешней сортировки, когда индекс строится частями на диске. Сортируется словарь ($V$ = 3,97 млн), а не все пары ($N$ = 364 млн), поэтому сортировка занимает 1,9~с из 148~с.
    \item Недостатки: нужен дополнительный буфер размером с массив; при каждом сравнении строки достаются по номеру терма из разных мест памяти (промахи кэша); весь словарь со списками должен помещаться в оперативную память. Для строк быстрее была бы поразрядная сортировка (MSD radix sort).
\end{itemize}
Альтернатива --- собрать все пары \enquote{терм, документ} и отсортировать их (BSBI). Она не требует хеш-таблицы и легко делается на диске, но сортирует в 92 раза больше элементов.

\subsection{Формат индекса}
Один файл \texttt{index.bin}; все числа --- беззнаковые, порядок байтов little-endian. Разделы идут подряд: заголовок (64 байта), таблица документов ($N \times 12$ байт), таблица термов ($T \times 16$ байт), списки документов ($P \times 4$ байт), пул строк.

\textbf{Заголовок} (64 байта):
\begin{center}
\small
\begin{tabular}{|r|r|l|p{0.5\textwidth}|}
\hline
Смещение & Байт & Поле & Значение \\
\hline
0 & 4 & \texttt{magic} & \texttt{BIDX} (\texttt{42 49 44 58}) \\
4 & 2 & \texttt{version} & версия формата, 1 \\
6 & 2 & \texttt{header\_size} & размер заголовка, 64 \\
8 & 4 & \texttt{flags} & флаги: бит 0 --- термы являются основами слов (стемминг, добавлен в ЛР4), бит 1 --- списки документов сжаты (добавлен в ЛР11), бит 2 --- в списках есть частоты термов, в записях документов --- длина вектора (добавлен в ЛР13), остальные биты зарезервированы, 0 \\
12 & 4 & \texttt{N} & число документов \\
16 & 4 & \texttt{T} & число термов \\
20 & 2 & \texttt{doc\_rec} & размер записи документа, 12 (16 при флаге 4, ЛР13) \\
22 & 2 & \texttt{term\_rec} & размер записи терма, 16 \\
24 & 8 & \texttt{P} & число записей в списках документов (пар терм--документ) \\
32 & 8 & \texttt{docs\_off} & смещение таблицы документов от начала файла \\
40 & 8 & \texttt{terms\_off} & смещение таблицы термов \\
48 & 8 & \texttt{postings\_off} & смещение списков документов \\
56 & 8 & \texttt{strings\_off} & смещение пула строк (до конца файла) \\
\hline
\end{tabular}
\end{center}

\textbf{Запись документа} (прямой индекс, 12 байт; запись с номером $i$ --- документ с docID $i$):
\begin{center}
\small
\begin{tabular}{|r|r|l|p{0.5\textwidth}|}
\hline
Смещение & Байт & Поле & Значение \\
\hline
0 & 4 & \texttt{title} & смещение заголовка в пуле строк \\
4 & 4 & \texttt{url} & смещение URL в пуле строк \\
8 & 4 & \texttt{tokens} & число токенов в документе \\
12 & 4 & \texttt{norm} & только при флаге 4 (ЛР13): длина вектора документа, \texttt{float} \\
\hline
\end{tabular}
\end{center}

\textbf{Запись терма} (16 байт; записи упорядочены по байтам строки терма):
\begin{center}
\small
\begin{tabular}{|r|r|l|p{0.5\textwidth}|}
\hline
Смещение & Байт & Поле & Значение \\
\hline
0 & 4 & \texttt{term} & смещение строки терма в пуле строк \\
4 & 4 & \texttt{df} & число документов с термом \\
8 & 8 & \texttt{first} & номер первой записи списка терма в разделе списков (при флаге 2 --- смещение списка в байтах, ЛР11) \\
\hline
\end{tabular}
\end{center}

\textbf{Списки документов}: $P$ чисел по 4 байта --- docID. Список терма --- записи с номерами от \texttt{first} до $\texttt{first} + \texttt{df} - 1$ (байты $\texttt{postings\_off} + 4 \cdot \texttt{first}$ и далее), docID по возрастанию, без повторов; списки идут в порядке термов. При флаге 2 (ЛР11) вместо чисел --- коды VByte разностей docID, при флаге 4 (ЛР13) каждая запись --- пара (docID, частота терма в документе).

\textbf{Пул строк}: строки подряд, каждая --- 2 байта длины $L$ и $L$ байтов UTF-8 без завершающего нуля. Сначала заголовок и URL каждого документа, затем строки термов.

\textbf{Как читается.} Терм ищется двоичным поиском по таблице термов: записи одного размера, поэтому $k$-я запись находится сразу (для 3,97 млн термов --- 22 сравнения). Список документов и запись прямого индекса тоже читаются по смещению, без разбора остального файла.

\textbf{Расширение формата.}
\begin{itemize}
    \item \texttt{version} --- номер версии; читатель отказывается открывать неизвестную версию.
    \item \texttt{header\_size}, \texttt{doc\_rec}, \texttt{term\_rec}: новые поля дописываются в конец заголовка и записей (например, длина документа в байтах, смещение позиций терма), а старый читатель шагает по записям с размером из заголовка и лишние поля пропускает.
    \item \texttt{flags} --- биты для новых возможностей (например, \enquote{в списках есть частоты термов}, \enquote{списки сжаты}).
    \item Новые разделы (координаты для цитатного поиска, зоны) дописываются после существующих, их смещения --- в новые поля заголовка.
\end{itemize}
